% \chapter{Theoretical Background}
\chapter{Kiến thức nền tảng}
% \hline
% \textit{In Chapter \thechapter, it presents about preliminary knowledge in this project.}
\textit{Trong chương này, chúng tôi tập trung xây dựng cơ sở lý thuyết đa ngành, tạo tiền đề toán học vững chắc cho việc phát triển và triển khai phương pháp đề xuất. Nội dung bao hàm việc hệ thống hóa các khái niệm về giải tích lồi, nón phối cảnh và các bài toán tối ưu hóa nón bậc hai (SOCP) vốn là nền tảng để chuyển đổi các ràng buộc phi lồi thành dạng lồi giải được hiệu quả. Đồng thời, các thuộc tính của đường cong Bézier, lý thuyết dòng mạng trong bài toán đường đi ngắn nhất và các kỹ thuật phân hoạch không gian an toàn cũng được trình bày chi tiết.}
\input{theoretical-background/convex-optimization.tex}
% % ref: Discrete Mathematics and Its Applications by Kenneth Rosen
% % ref: Combinatorial Optimization: Theory and Algorithms
% \section{Graph Theory and the Shortest Path Problem}

% \subsection{Graph Theory Fundamentals}

% \subsubsection{Basic Definitions and Concepts}

% Formally, a graph is defined as an ordered pair $G = (\mathcal{V}, \mathcal{E})$, comprising a set of vertices $\mathcal{V}$ and a set of edges $\mathcal{E}$, where each edge connects a pair of vertices. The structural classification of the graph depends on the ordering of these pairs: unordered pairs $\{u, w\}$ denote an undirected graph, whereas ordered pairs $(u, w)$ characterize a directed graph. In the context of optimization, the graph is typically augmented with a cost function $c: \mathcal{E} \rightarrow \mathbb{R}$, assigning a real-valued weight to each edge.

% Key topological properties include the vertex degree, which counts incident edges; in directed graphs, this is distinguished into in-degree ($\text{deg}^-$) and out-degree ($\text{deg}^+$). Connectivity is described via paths—finite sequences of distinct vertices and edges—where the path length is the aggregate cost of its constituent edges. A cycle is defined as a path that originates and terminates at the same vertex. Optimization problems generally operate on a search domain defined by subgraphs $H = (W, F)$, formed by subsets of vertices $W \subseteq \mathcal{V}$ and edges $F \subseteq \mathcal{E}$ that maintain incidence relationships.

% \subsubsection{Combinatorial Optimization on Graphs}

% Combinatorial optimization problems on graphs generally seek a substructure that minimizes an objective function subject to validity constraints. For a weighted graph $G = (\mathcal{V}, \mathcal{E})$ and a set of feasible subgraphs $\mathcal{H}$, the problem is to identify a subgraph $H = (W, F) \in \mathcal{H}$ that minimizes the cumulative weight of its components. This is formally expressed as:

% \begin{align}
% \text{minimize} \quad & \sum_{v \in W} c_v + \sum_{e \in F} c_e \tag{1.1a} \\
% \text{subject to} \quad & H=(W, F) \in \mathcal{H}. \tag{1.1b}
% \end{align}

% To utilize standard mathematical programming techniques, such as Branch and Bound, the problem is transformed into an Integer Linear Programming (ILP) formulation. This involves parameterizing the search space using incidence vectors $\mathbf{y}^H \in \{0, 1\}^{|\mathcal{V} \cup \mathcal{E}|}$, where binary variables indicate the inclusion of specific vertices or edges in the solution. Geometrically, the set of valid subgraphs corresponds to a polytope $\mathcal{Y}$ in Euclidean space. The intersection of this polytope with the integer grid, $\mathcal{Y} \cap \{0, 1\}^{|\mathcal{V} \cup \mathcal{E}|}$, yields the incidence vectors of feasible subgraphs. Consequently, the discrete optimization problem is recast as a linear optimization over an integer domain:

% \begin{align}
% \text{minimize} \quad & \sum_{v \in \mathcal{V}} c_v y_v + \sum_{e \in \mathcal{E}} c_e y_e \tag{1.2a} \\
% \text{subject to} \quad & \mathbf{y} \in \mathcal{Y} \cap \{0, 1\}^{|\mathcal{V} \cup \mathcal{E}|}. \tag{1.2b}
% \end{align}

% This formulation allows for the application of convex analysis, ensuring that the optimal solution to the algebraic model corresponds exactly to the optimal substructure in the original graph problem.

% \subsection{The Shortest Path Problem (SPP)}

% \subsubsection{Definition and Classical Algorithms}

% The Shortest Path Problem (SPP) on a weighted graph $G=(\mathcal{V}, \mathcal{E})$ involves identifying a path $P$ between a source $s$ and a destination $t$ such that the summation of edge weights is minimized. If $P$ comprises a sequence of edges $e_1, \ldots, e_k$, the objective is to minimize $\sum_{i=1}^{k} c_{e_i}$.

% Classical approaches to solving the SPP vary based on edge weight characteristics. Dijkstra's algorithm employs a greedy strategy suitable for graphs with non-negative weights ($c_e \ge 0$). For graphs containing negative edge weights, provided no negative cycles exist, the Bellman-Ford algorithm is utilized, albeit with higher computational complexity. The Floyd-Warshall algorithm extends this to the all-pairs shortest path problem, accommodating negative edges but remaining sensitive to negative cycles.

% \subsubsection{Network Flow Formulation for the Shortest Path Problem}

% The Shortest Path Problem on a directed graph $G = (\mathcal{V}, \mathcal{E})$ with non-negative edge costs $c_{uv} \ge 0$ can be rigorously modeled as a Minimum Cost Flow problem involving a single unit of flow. We define binary decision variables $x_{uv} \in \{0, 1\}$ for each edge $(u, v) \in \mathcal{E}$, where $x_{uv} = 1$ indicates the edge is part of the optimal path. The optimization model is formulated as follows:

% \begin{subequations}
% \begin{align}
% \text{minimize} \quad & \sum_{(u, v) \in \mathcal{E}} c_{uv} x_{uv} \label{eq:spp_obj} \\
% \text{subject to} \quad & \sum_{v \in \mathcal{V}^+(u)} x_{uv} - \sum_{v \in \mathcal{V}^-(u)} x_{vu} =
% \begin{cases}
% 1 & \text{if } u = s \\
% -1 & \text{if } u = t \\
% 0 & \text{otherwise}
% \end{cases}, \quad \forall u \in \mathcal{V} \label{eq:flow_conservation} \\
% & \sum_{v \in \mathcal{V}^-(u)} x_{vu} \le 1, \quad \forall u \in \mathcal{V} \setminus \{s, t\} \label{eq:node_capacity} \\
% & x_{uv} \ge 0, \quad \forall (u, v) \in \mathcal{E} \label{eq:non_negativity}
% \end{align}
% \end{subequations}

% Here, $\mathcal{V}^+(u)$ and $\mathcal{V}^-(u)$ denote the sets of successors and predecessors of node $u$, respectively. The constraints in this formulation enforce specific structural properties. Equation \eqref{eq:flow_conservation} represents flow conservation (Kirchhoff's law), ensuring that a net flow of one unit leaves the source $s$, enters the target $t$, and is conserved at all intermediate nodes. Equation \eqref{eq:node_capacity} imposes a node capacity constraint, restricting the inflow at any intermediate vertex to unity, thereby enforcing a simple path structure where no vertex is revisited. Given non-negative costs, cycles are inherently suboptimal, rendering explicit cycle-elimination constraints unnecessary.

% A significant advantage of this formulation lies in the algebraic properties of the constraint matrix derived from \eqref{eq:flow_conservation}. This node-arc incidence matrix exhibits Total Unimodularity (TUM). According to combinatorial optimization theory, if the constraint matrix is TUM and the right-hand side vector is integral, every basic feasible solution of the Linear Programming (LP) relaxation is guaranteed to be integral~\cite{Ahuja1993}. Consequently, despite the continuous non-negativity constraints in \eqref{eq:non_negativity}, the problem can be solved efficiently using standard LP algorithms to yield a valid binary path, obviating the need for computationally intensive integer programming methods.



\section{Lý thuyết Đồ thị và Bài toán Đường đi Ngắn nhất}

\subsection{Cơ sở Lý thuyết Đồ thị}

\subsubsection{Các định nghĩa và khái niệm cơ bản}

Một cách hình thức, một đồ thị được định nghĩa là một cặp có thứ tự
$G = (\mathcal{V}, \mathcal{E})$, bao gồm một tập các đỉnh $\mathcal{V}$ và một
tập các cạnh $\mathcal{E}$, trong đó mỗi cạnh nối một cặp đỉnh. Việc phân loại
cấu trúc của đồ thị phụ thuộc vào tính có thứ tự của các cặp này: các cặp không
có thứ tự $\{u, w\}$ biểu diễn đồ thị vô hướng, trong khi các cặp có thứ tự
$(u, w)$ đặc trưng cho đồ thị có hướng. Trong bối cảnh tối ưu hóa, đồ thị thường
được mở rộng bằng một hàm chi phí
$c: \mathcal{E} \rightarrow \mathbb{R}$, gán một trọng số thực cho mỗi cạnh.

Các tính chất tô pô quan trọng bao gồm bậc của đỉnh, là số cạnh liên thuộc với
đỉnh đó; trong đồ thị có hướng, khái niệm này được phân biệt thành bậc vào
($\text{deg}^-$) và bậc ra ($\text{deg}^+$). Tính liên thông được mô tả thông qua
các đường đi—là các dãy hữu hạn gồm các đỉnh và cạnh phân biệt—trong đó độ dài
đường đi được xác định bằng tổng chi phí của các cạnh cấu thành. Một chu trình
được định nghĩa là một đường đi bắt đầu và kết thúc tại cùng một đỉnh. Các bài
toán tối ưu thường được xét trên một miền tìm kiếm được xác định bởi các đồ thị
con $H = (W, F)$, hình thành từ các tập con của đỉnh $W \subseteq \mathcal{V}$
và các cạnh $F \subseteq \mathcal{E}$, đồng thời bảo toàn các quan hệ liên thuộc.

\subsubsection{Tối ưu hóa tổ hợp trên đồ thị}

Các bài toán tối ưu hóa tổ hợp trên đồ thị thường nhằm tìm một cấu trúc con sao
cho hàm mục tiêu đạt giá trị nhỏ nhất, đồng thời thỏa mãn các ràng buộc hợp lệ.
Với một đồ thị có trọng số $G = (\mathcal{V}, \mathcal{E})$ và một tập các đồ thị
con khả thi $\mathcal{H}$, bài toán là xác định một đồ thị con
$H = (W, F) \in \mathcal{H}$ sao cho tổng trọng số của các thành phần của nó là
nhỏ nhất. Bài toán này được biểu diễn hình thức như sau:
\begin{align}
\text{tối thiểu hóa} \quad & \sum_{v \in W} c_v + \sum_{e \in F} c_e \tag{1.1a} \\
\text{với điều kiện} \quad & H=(W, F) \in \mathcal{H}. \tag{1.1b}
\end{align}

Để có thể áp dụng các kỹ thuật lập trình toán học tiêu chuẩn, chẳng hạn như
Branch-and-Bound, bài toán được chuyển đổi sang dạng Quy hoạch Tuyến tính
Nguyên (Integer Linear Programming -- ILP). Quá trình này bao gồm việc tham số
hóa không gian tìm kiếm bằng các vectơ liên thuộc
$\mathbf{y}^H \in \{0,1\}^{|\mathcal{V} \cup \mathcal{E}|}$, trong đó các biến nhị
phân biểu thị việc một đỉnh hoặc một cạnh cụ thể có được chọn vào nghiệm hay
không. Về mặt hình học, tập các đồ thị con hợp lệ tương ứng với một đa diện
$\mathcal{Y}$ trong không gian Euclid. Giao của đa diện này với lưới số nguyên,
$\mathcal{Y} \cap \{0,1\}^{|\mathcal{V} \cup \mathcal{E}|}$, tạo thành các vectơ
liên thuộc của các đồ thị con khả thi. Do đó, bài toán tối ưu rời rạc ban đầu
được tái biểu diễn thành một bài toán tối ưu tuyến tính trên miền nguyên:
\begin{align}
\text{tối thiểu hóa} \quad &
\sum_{v \in \mathcal{V}} c_v y_v + \sum_{e \in \mathcal{E}} c_e y_e \tag{1.2a} \\
\text{với điều kiện} \quad &
\mathbf{y} \in \mathcal{Y} \cap \{0,1\}^{|\mathcal{V} \cup \mathcal{E}|}. \tag{1.2b}
\end{align}

Cách phát biểu này cho phép áp dụng các công cụ của phân tích lồi, đồng thời đảm
bảo rằng nghiệm tối ưu của mô hình đại số tương ứng chính xác với cấu trúc con
tối ưu của bài toán đồ thị ban đầu.

\subsection{Bài toán Đường đi Ngắn nhất (SPP)}

\subsubsection{Định nghĩa và các thuật toán cổ điển}

Bài toán Đường đi Ngắn nhất (Shortest Path Problem -- SPP) trên một đồ thị có
trọng số $G = (\mathcal{V}, \mathcal{E})$ nhằm xác định một đường đi $P$ từ đỉnh
nguồn $s$ đến đỉnh đích $t$ sao cho tổng trọng số các cạnh là nhỏ nhất. Nếu $P$
bao gồm một dãy các cạnh $e_1, \ldots, e_k$, mục tiêu là tối thiểu hóa
$\sum_{i=1}^{k} c_{e_i}$.

Các phương pháp cổ điển để giải SPP phụ thuộc vào đặc tính của trọng số cạnh.
Thuật toán Dijkstra sử dụng chiến lược tham lam và phù hợp với các đồ thị có
trọng số không âm ($c_e \ge 0$). Đối với các đồ thị có cạnh mang trọng số âm,
miễn là không tồn tại chu trình âm, thuật toán Bellman--Ford được sử dụng, mặc
dù có độ phức tạp tính toán cao hơn. Thuật toán Floyd--Warshall mở rộng bài toán
này cho trường hợp tìm đường đi ngắn nhất giữa mọi cặp đỉnh, cho phép tồn tại
trọng số âm nhưng vẫn nhạy cảm với sự xuất hiện của các chu trình âm.

\subsubsection{Mô hình dòng mạng cho Bài toán Đường đi Ngắn nhất}

Bài toán Đường đi Ngắn nhất trên một đồ thị có hướng
$G = (\mathcal{V}, \mathcal{E})$ với chi phí cạnh không âm $c_{uv} \ge 0$ có thể
được mô hình hóa một cách chặt chẽ như một bài toán Dòng Chi phí Tối thiểu với
một đơn vị dòng. Ta định nghĩa các biến quyết định nhị phân
$x_{uv} \in \{0,1\}$ cho mỗi cạnh $(u,v) \in \mathcal{E}$, trong đó
$x_{uv} = 1$ cho biết cạnh đó thuộc đường đi tối ưu. Mô hình tối ưu được phát
biểu như sau:
\begin{subequations}
\begin{align}
\text{minimize} \quad & \sum_{(u, v) \in \mathcal{E}} c_{uv} x_{uv} \label{eq:spp_obj} \\
\text{subject to} \quad & \sum_{v \in \mathcal{V}^+(u)} x_{uv} - \sum_{v \in \mathcal{V}^-(u)} x_{vu} =
\begin{cases}
1 & \text{if } u = s \\
-1 & \text{if } u = t \\
0 & \text{otherwise}
\end{cases}, \quad \forall u \in \mathcal{V} \label{eq:flow_conservation} \\
& \sum_{v \in \mathcal{V}^-(u)} x_{vu} \le 1, \quad \forall u \in \mathcal{V} \setminus \{s, t\} \label{eq:node_capacity} \\
& x_{uv} \ge 0, \quad \forall (u, v) \in \mathcal{E} \label{eq:non_negativity}
\end{align}
\end{subequations}

Ở đây, $\mathcal{V}^+(u)$ và $\mathcal{V}^-(u)$ lần lượt ký hiệu tập các đỉnh kế
tiếp (successors) và các đỉnh tiền nhiệm (predecessors) của đỉnh $u$. Các ràng
buộc trong mô hình này đảm bảo những tính chất cấu trúc cụ thể. Phương trình
\eqref{eq:flow_conservation} biểu diễn điều kiện bảo toàn dòng (định luật
Kirchhoff), đảm bảo rằng một đơn vị dòng rời khỏi đỉnh nguồn $s$, đi vào đỉnh
đích $t$, và được bảo toàn tại tất cả các đỉnh trung gian. Ràng buộc
\eqref{eq:node_capacity} áp đặt giới hạn dung lượng nút, giới hạn dòng đi vào
mỗi đỉnh trung gian không vượt quá một đơn vị, qua đó đảm bảo cấu trúc đường đi
đơn, trong đó không có đỉnh nào bị ghé thăm nhiều hơn một lần. Với chi phí
không âm, các chu trình luôn là không tối ưu, do đó không cần thiết phải bổ sung
các ràng buộc loại bỏ chu trình một cách tường minh.

Một ưu điểm quan trọng của cách mô hình hóa này nằm ở các tính chất đại số của
ma trận ràng buộc thu được từ~\eqref{eq:flow_conservation}. Ma trận liên thuộc
đỉnh-cạnh này có tính \textit{hoàn toàn đơn mô} (Total Unimodularity--TUM).
Theo lý thuyết tối ưu hóa tổ hợp, nếu ma trận ràng buộc là TUM và vectơ vế phải
là nguyên, thì mọi nghiệm khả thi cơ sở của bài toán thư giãn Quy hoạch Tuyến
tính (LP) đều được đảm bảo là nguyên~\cite{AhujaMagnantiOrlin1993}. Do đó, mặc dù chỉ áp đặt
các ràng buộc không âm liên tục trong~\eqref{eq:non_negativity}, bài toán vẫn có
thể được giải hiệu quả bằng các thuật toán LP tiêu chuẩn để thu được một đường
đi nhị phân hợp lệ, mà không cần đến các phương pháp lập trình nguyên có chi phí
tính toán cao.

\section{Sinh Quỹ đạo Kinodynamic}\label{sec:kinodynamic_theory}

Lập kế hoạch chuyển động hình học (geometric motion planning) tìm kiếm một đường đi
không va chạm trong không gian cấu hình của robot mà không xét đến các đặc tính vật lý
của hệ thống.
Tuy nhiên, trong thực tế, robot là một hệ cơ học mà trạng thái của nó chỉ có thể tiến triển
theo thời gian theo các phương trình vi phân xác định được dẫn dắt bởi các đầu vào cơ cấu
chấp hành.
Một đường đi không va chạm về mặt hình học vẫn có thể không khả thi về mặt vật lý nếu nó
đòi hỏi, chẳng hạn, sự thay đổi tức thời của vận tốc hoặc các lực vượt quá giới hạn bão hòa
của cơ cấu chấp hành.
Lập kế hoạch chuyển động kinodynamic mở rộng bài toán hình học bằng cách yêu cầu quỹ đạo
được lập kế hoạch phải đồng thời tránh vật cản, thỏa mãn các ràng buộc vi phân do động lực
học hệ thống áp đặt, và tối ưu hóa một tiêu chí hiệu năng do người dùng chỉ định.

Mục này xây dựng bài toán lập kế hoạch kinodynamic dưới dạng toán học chặt chẽ
(Mục~\ref{subsec:kinodynamic_problem}), giới thiệu tham số hóa quỹ đạo đa thức bằng
đường cong Bézier như một kỹ thuật để đưa bài toán về dạng hữu hạn chiều
(Mục~\ref{subsec:bezier_parameterization}), và trình bày các ràng buộc liên tục đảm bảo
chuyển động trơn qua các điểm chuyển tiếp giữa các đoạn quỹ đạo
(Mục~\ref{subsec:continuity_constraints}).

\subsection{Bài toán Lập kế hoạch Kinodynamic}\label{subsec:kinodynamic_problem}

Xét một hệ thống robot có trạng thái $x(t) \in \mathcal{X} \subseteq \mathbb{R}^{n_x}$
và đầu vào điều khiển $u(t) \in \mathcal{U} \subseteq \mathbb{R}^{n_u}$ tiến triển theo
phương trình vi phân thường
\begin{equation}
    \dot{x}(t) = f\!\left(x(t),\, u(t)\right), \quad t \in [0, T],
    \label{eq:kinodynamic_dynamics}
\end{equation}
trong đó $f \colon \mathbb{R}^{n_x} \times \mathbb{R}^{n_u} \to \mathbb{R}^{n_x}$ là ánh xạ
động lực học của hệ và $T > 0$ là độ dài chân trời lập kế hoạch, có thể được cố định trước
hoặc được coi là biến tối ưu.
Vector trạng thái thường mã hóa các vị trí và vận tốc suy rộng của robot, và tùy theo
mô hình động lực học có thể bao gồm thêm các đạo hàm bậc cao hơn.

Một quỹ đạo kinodynamic là một cặp $\left(x(\cdot), u(\cdot)\right)$ xác định trên $[0,T]$
thỏa mãn phương trình động lực học~\eqref{eq:kinodynamic_dynamics} cùng với các yêu cầu sau.

\begin{enumerate}
    \item \textbf{Điều kiện biên.}
          Quỹ đạo phải xuất phát từ trạng thái đầu được chỉ định và đến được vùng đích mong muốn:
          \begin{equation}
              x(0) = x_{\mathrm{start}}, \qquad x(T) \in \mathcal{X}_{\mathrm{goal}},
              \label{eq:boundary_conditions}
          \end{equation}
          trong đó $\mathcal{X}_{\mathrm{goal}} \subseteq \mathcal{X}$ là tập đích.

    \item \textbf{Ràng buộc an toàn.}
          Quỹ đạo trạng thái phải nằm trong vùng không có vật cản
          $\mathcal{X}_{\mathrm{free}} \subseteq \mathcal{X}$ tại mọi thời điểm:
          \begin{equation}
              x(t) \in \mathcal{X}_{\mathrm{free}}, \quad \forall\, t \in [0, T].
              \label{eq:safety_constraint}
          \end{equation}

    \item \textbf{Giới hạn cơ cấu chấp hành và động học.}
          Đầu vào điều khiển phải nằm trong tập khả thi của cơ cấu chấp hành:
          \begin{equation}
              u(t) \in \mathcal{U}, \quad \forall\, t \in [0, T].
              \label{eq:control_constraint}
          \end{equation}
          Trong nhiều bài toán thực tế, $\mathcal{U}$ được đặc trưng bởi các giới hạn tường minh
          lên vận tốc, gia tốc và độ giật dọc theo quỹ đạo,
          ví dụ\ $\|\dot{x}(t)\| \le v_{\max}$, $\|\ddot{x}(t)\| \le a_{\max}$
\end{enumerate}

Mục tiêu là tìm một quỹ đạo tối thiểu hóa phiếm hàm chi phí có dạng
\begin{equation}
    J\!\left(x(\cdot), u(\cdot)\right)
    = \int_{0}^{T} L\!\left(x(t), u(t)\right) \mathrm{d}t + E\!\left(x(T)\right),
    \label{eq:kinodynamic_cost}
\end{equation}
trong đó $L$ là running cost phạt, chẳng hạn, nỗ lực điều khiển hoặc độ lệch so với
quỹ đạo tham chiếu, và $E$ là terminal cost.

Có hai khó khăn cấu trúc phân biệt bài toán này với bài toán hình học thuần túy.
Thứ nhất, ràng buộc an toàn~\eqref{eq:safety_constraint} phải được thỏa mãn liên tục
trên toàn bộ khoảng $[0, T]$, không chỉ tại một tập hữu hạn các điểm rời rạc,
và $\mathcal{X}_{\mathrm{free}}$ nói chung là phi lồi do sự hiện diện của vật cản.
Thứ hai, phiếm hàm chi phí~\eqref{eq:kinodynamic_cost} và ràng buộc động lực
học~\eqref{eq:kinodynamic_dynamics} có tính chất vô hạn chiều, vì các biến quyết định
$x(\cdot)$ và $u(\cdot)$ là các hàm theo thời gian.
Để giải bài toán này bằng phương pháp số, cần phải đưa nó về dạng bài toán hữu hạn chiều
bằng cách chọn một tham số hóa quỹ đạo phù hợp.

\subsection{Tham số hóa Quỹ đạo bằng Đường cong Bézier}\label{subsec:bezier_parameterization}

Một hướng tiếp cận tự nhiên để tham số hóa hữu hạn chiều là biểu diễn quỹ đạo dưới dạng
đa thức với các hệ số đóng vai trò là biến quyết định.
Trong số các biểu diễn đa thức, đường cong Bézier đặc biệt phù hợp với lập kế hoạch
kinodynamic vì các tính chất hình học và giải tích của chúng cho phép áp đặt các ràng buộc
an toàn, độ trơn và tính khả thi động học trực tiếp dưới dạng các điều kiện đại số trên
các hệ số.

\subsubsection{Định nghĩa}

Một đường cong Bézier bậc $d$ trên khoảng tham số $[0,1]$ được xác định bởi
$d+1$ \emph{điểm điều khiển} $\gamma_0, \ldots, \gamma_d \in \mathbb{R}^{n}$
thông qua các đa thức cơ sở Bernstein.
Đa thức Bernstein thứ $k$ bậc $d$ được định nghĩa là
\begin{equation}
    \beta_{k,d}(s) := \binom{d}{k} s^{k}(1-s)^{d-k},
    \quad k = 0, \ldots, d, \quad s \in [0,1].
    \label{eq:bernstein}
\end{equation}
Theo định lý nhị thức, các đa thức này không âm và phân hoạch đơn vị:
$\sum_{k=0}^{d} \beta_{k,d}(s) = 1$ với mọi $s \in [0,1]$.
Do đó, với mỗi $s$ cố định, các giá trị $\{\beta_{k,d}(s)\}_{k=0}^{d}$ tạo thành một
tập hệ số tổ hợp lồi.
Đường cong Bézier được định nghĩa là trung bình có trọng số tương ứng của các điểm điều khiển:
\begin{equation}
    \gamma(s) := \sum_{k=0}^{d} \beta_{k,d}(s)\,\gamma_k,
    \quad s \in [0,1].
    \label{eq:bezier_curve}
\end{equation}

\subsubsection{Các tính chất liên quan đến lập kế hoạch quỹ đạo}

Ba tính chất của đường cong Bézier được khai thác trực tiếp trong lập kế hoạch quỹ đạo
kinodynamic.

\textbf{Nội suy điểm đầu cuối.}
Đường cong đi qua chính xác điểm điều khiển đầu tiên và điểm điều khiển cuối cùng:
\begin{equation}
    \gamma(0) = \gamma_0, \qquad \gamma(1) = \gamma_d.
    \label{eq:endpoint_interpolation}
\end{equation}
Tính chất này cho phép áp đặt các điều kiện biên về vị trí và, thông qua công thức đạo hàm
bên dưới, cả về vận tốc, dưới dạng các ràng buộc đẳng thức tuyến tính trên các điểm điều khiển.

\textbf{Tính chất bao lồi.}
Toàn bộ đường cong nằm bên trong bao lồi của các điểm điều khiển của nó:
\begin{equation}
    \gamma(s) \in \operatorname{conv}\!\left(\{\gamma_k\}_{k=0}^{d}\right),
    \quad \forall\, s \in [0,1].
    \label{eq:convex_hull}
\end{equation}
Suy ra rằng nếu tất cả $d+1$ điểm điều khiển nằm trong một tập lồi $\mathcal{S}$,
thì $\gamma(s) \in \mathcal{S}$ với mọi $s$.
Nguyên tắc bao hàm này chuyển đổi ràng buộc an toàn liên tục~\eqref{eq:safety_constraint}
thành một tập hữu hạn các ràng buộc thuộc tập trên các điểm điều khiển, với điều kiện vùng
an toàn có thể được xấp xỉ hoặc phân hoạch thành các tập lồi.

\textbf{Công thức đạo hàm (Hodograph).}
Đạo hàm của $\gamma$ theo tham số $s$ cũng là một đường cong Bézier bậc $d-1$,
với các điểm điều khiển được tạo thành từ các sai phân hữu hạn bậc một của các điểm
điều khiển ban đầu:
\begin{equation}
    \dot{\gamma}(s) = \sum_{k=0}^{d-1} \beta_{k,d-1}(s)\,\dot{\gamma}_k,
    \qquad \dot{\gamma}_k = d\!\left(\gamma_{k+1} - \gamma_k\right).
    \label{eq:bezier_derivative}
\end{equation}
Áp dụng đệ quy~\eqref{eq:bezier_derivative} cho phép tính gia tốc và các đạo hàm bậc cao
hơn dưới dạng đường cong Bézier có bậc giảm dần, với các điểm điều khiển được biểu diễn
dưới dạng tổ hợp tuyến tính của các $\gamma_k$ ban đầu.
Do đó, các ràng buộc lên vận tốc và gia tốc dọc theo quỹ đạo --- chẳng hạn các giới hạn
động học trong~\eqref{eq:control_constraint} --- quy về các ràng buộc bất đẳng thức tuyến
tính trên các sai phân điểm điều khiển, và có thể được áp đặt đồng thời với tính chất bao lồi.

\subsection{Ràng buộc Độ trơn tại Điểm chuyển tiếp}\label{subsec:continuity_constraints}

Với các quỹ đạo trải dài qua nhiều vùng không gian làm việc có tính chất động học khác nhau,
một đoạn Bézier duy nhất bậc vừa phải có thể không đủ linh hoạt.
Biện pháp thực tế là sử dụng quỹ đạo \emph{Bézier từng đoạn} (piecewise Bézier) gồm $M$
đoạn ghép nối đầu-đuôi.
Gọi đoạn thứ $i$ là đường cong Bézier bậc $d$ được xác định bởi các điểm điều khiển
$a^{(i)}_0, \ldots, a^{(i)}_d$, với $i = 1, \ldots, M$.
Việc ghép nối các đoạn đảm bảo đường đi liên thông, nhưng cần thêm các điều kiện đại số
tại mỗi điểm nối để đảm bảo tính khả vi của quỹ đạo tổng thể.

\textbf{Liên tục $C^0$ (vị trí).}
Điểm cuối của đoạn $i$ trùng với điểm đầu của đoạn $i+1$:
\begin{equation}
    a^{(i)}_d = a^{(i+1)}_0.
    \label{eq:c0_continuity}
\end{equation}

\textbf{Liên tục $C^1$ (vận tốc).}
Vận tốc ra của đoạn $i$ khớp với vận tốc vào của đoạn $i+1$,
điều này theo~\eqref{eq:bezier_derivative} được dịch thành đẳng thức của các sai phân bậc
một tại biên:
\begin{equation}
    a^{(i)}_d - a^{(i)}_{d-1} = a^{(i+1)}_1 - a^{(i+1)}_0.
    \label{eq:c1_continuity}
\end{equation}

\textbf{Liên tục $C^2$ (gia tốc).}
Gia tốc ra của đoạn $i$ khớp với gia tốc vào của đoạn $i+1$,
được biểu diễn qua các sai phân bậc hai tại biên:
\begin{equation}
    a^{(i)}_d - 2a^{(i)}_{d-1} + a^{(i)}_{d-2}
    = a^{(i+1)}_2 - 2a^{(i+1)}_1 + a^{(i+1)}_0.
    \label{eq:c2_continuity}
\end{equation}

Các điều kiện~\eqref{eq:c0_continuity}--\eqref{eq:c2_continuity} là các ràng buộc đẳng thức
tuyến tính trên các điểm điều khiển và do đó bảo toàn cấu trúc đại số của bài toán tối ưu.
Việc yêu cầu liên tục $C^1$ hay $C^2$ phụ thuộc vào bậc của mô hình động lực học robot:
mô hình tích phân đôi (double-integrator) chỉ cần $C^1$, trong khi mô hình có đầu vào mô-men
xoắn tường minh thường đòi hỏi $C^2$ để tránh các lệnh lực xung tại điểm nối giữa các đoạn.

Sau khi đặt tham số hóa Bézier từng đoạn, bài toán lập kế hoạch kinodynamic quy về bài toán
tối ưu hữu hạn chiều: các biến quyết định là các điểm điều khiển $\{a^{(i)}_k\}$ và, khi áp
dụng, chân trời $T$; hàm mục tiêu là phiếm hàm chi phí~\eqref{eq:kinodynamic_cost} được biểu
diễn qua các điểm điều khiển; và các ràng buộc bao gồm điều kiện
biên~\eqref{eq:boundary_conditions}, các ràng buộc an toàn suy từ tính chất bao
lồi~\eqref{eq:convex_hull}, các giới hạn động học được áp đặt qua công thức đạo
hàm~\eqref{eq:bezier_derivative}, và các điều kiện tại điểm
nối~\eqref{eq:c0_continuity}--\eqref{eq:c2_continuity}.

Cấu trúc này --- hàm mục tiêu tích phân, ràng buộc vi phân trên quỹ đạo, điều kiện biên cố
định, và các ràng buộc bất đẳng thức trên đường --- chính xác là cấu trúc của một Bài toán
Điều khiển Tối ưu (Optimal Control Problem, OCP).
Mục tiếp theo hệ thống hóa mối liên hệ này và phát triển khung OCP làm nền tảng cho phương
pháp tối ưu hóa quỹ đạo được áp dụng trong đồ án này.

% \section{Convex Decomposition}
\section{Phân hoạch Lồi}
\label{sec:convex-decomposition}

Phần này cung cấp tổng quan về phân hoạch lồi, một kỹ thuật cơ bản được sử dụng rộng rãi trong
hình học tính toán, mô phỏng vật lý, và đặc biệt là lập kế hoạch chuyển động robot.

\subsection{Định nghĩa}

\textbf{Phân hoạch lồi (Convex Decomposition)} là quá trình chia một hình học phức tạp (ví dụ:
đa giác hoặc đa diện không lồi) thành một tập hợp các \textbf{vùng con lồi} sao cho hợp của các
vùng này khôi phục lại toàn bộ hình ban đầu và các vùng con không chồng lấn nhau, ngoại trừ biên
chung.

Về mặt hình thức, với một miền hình học $P \subset \mathbb{R}^n$, một \textit{phân hoạch lồi} của
$P$ là tập hợp $\{P_1, P_2, \dots, P_k\}$ sao cho:
\[
    P = \bigcup_{i=1}^{k} P_i,
    \qquad
    P_i \cap P_j = \partial P_i \cap \partial P_j \text{ với mọi } i \ne j,
\]
và mỗi $P_i$ là một tập lồi.
Mục tiêu là tìm phân hoạch này sao cho tối thiểu hóa số lượng các vùng lồi, giảm thiểu phần chồng
chéo, hoặc tối ưu chất lượng hình học của các vùng lồi được tạo ra (tránh các vùng quá ``mảnh'' hoặc
quá ``dài'').

\subsection{Phân loại Phương pháp Phân hoạch Lồi}
\label{sec:convex-decomposition-types}

Trong lĩnh vực hình học tính toán và hoạch định chuyển động, các phương pháp phân hoạch lồi thường
được chia thành hai loại chính:
\textbf{Phân hoạch Lồi Chính xác (Exact Convex Decomposition~-- ECD)}
và
\textbf{Phân hoạch Lồi Xấp xỉ (Approximate Convex Decomposition~-- ACD)}.
Sự khác biệt giữa hai loại này phản ánh sự đánh đổi cơ bản giữa \textit{tính chính xác hình học}
và \textit{hiệu quả tính toán}.

\subsubsection{Phân hoạch Lồi Chính xác}

Phân hoạch Lồi Chính xác (ECD) là quá trình chia một hình dạng không lồi thành các \textbf{thành
phần con hoàn toàn lồi}, không chồng lấn, sao cho hợp của chúng tái tạo lại chính xác hình dạng ban
đầu.
Mục tiêu thông thường là tối thiểu hóa số lượng vùng lồi.

Tuy nhiên, ECD gặp hai thách thức lớn:
\begin{itemize}
    \item \textbf{Độ phức tạp tính toán cao:}
    Bài toán tìm phân hoạch lồi tối thiểu đã được chứng minh là NP-hard đối với đa giác có lỗ,
    nên không tồn tại thuật toán hiệu quả cho các trường hợp tổng quát.
    \item \textbf{Tính khả dụng hạn chế:}
    Ngay cả khi tính toán được, ECD thường tạo ra số lượng vùng rất lớn,
    khiến việc lưu trữ, xử lý và tối ưu hóa tiếp theo trở nên tốn kém.
\end{itemize}
Do đó, các phương pháp ECD chủ yếu mang tính lý thuyết và dùng làm chuẩn so sánh,
trong khi hầu hết các ứng dụng thực tế trong robot và mô phỏng đều dựa vào các phương pháp xấp xỉ.

\subsubsection{Phân hoạch Lồi Xấp xỉ}

Trước những giới hạn của ECD, các phương pháp \textbf{Phân hoạch Lồi Xấp xỉ (ACD)} được phát triển
để đạt sự cân bằng giữa độ chính xác và hiệu quả.
ACD cho phép các thành phần không hoàn toàn lồi, mà chỉ ``gần lồi'' trong một mức dung sai lõm
$\tau$ do người dùng xác định, hoặc chỉ bao phủ một phần không gian tự do.
Triết lý cốt lõi của ACD là chấp nhận một \textbf{mức độ lõm nhỏ có kiểm soát} để đổi lấy:
\begin{itemize}
    \item \textbf{Hiệu quả tính toán cao hơn:}
    Các thuật toán ACD, chẳng hạn như phương pháp của Lien và Amato~\cite{LienAmato2006},
    chạy nhanh hơn đáng kể so với ECD.
    \item \textbf{Giảm số lượng vùng:}
    Cho phép biểu diễn hình dạng với ít vùng lồi hơn, đơn giản hơn và dễ xử lý hơn.
    \item \textbf{Kiểm soát mức độ chi tiết:}
    Một số thuật toán ACD, chẳng hạn như phương pháp của Lien và Amato, cho phép người dùng điều
    chỉnh tham số lõm $\tau$ để cân bằng giữa độ chính xác và số lượng vùng lồi.
    Đối với phương pháp Visibility Clique Cover (VCC)~\cite{DeitsTedrake2023VCC}, người dùng có
    thể chọn chỉ số bao phủ thể hiện mức độ bao phủ của các vùng lồi so với không gian tự do ban đầu.
\end{itemize}

Do những ưu điểm trên, ACD được ưu tiên áp dụng trong hầu hết các hệ thống robot thực tế,
bao gồm cả khuôn khổ được đề xuất trong đồ án này.
Cụ thể, phương pháp ACD tạo ra một tập hợp các vùng lồi phủ không gian tự do, từ đó xây dựng đồ thị
$\mathcal{G}$ cho bài toán GCS.
Chất lượng của phân hoạch ảnh hưởng trực tiếp đến số lượng nút trong đồ thị và do đó đến độ phức
tạp của bài toán tối ưu hóa~-- một phân hoạch tốt cần tạo ra đủ ít vùng để giữ cho bài toán gọn
nhẹ, đồng thời đủ chi tiết để robot có thể đi qua các hành lang hẹp trong môi trường.

% 
\section{Tối ưu hóa Nguyên-Hỗn hợp}
\label{sec:mip_background}

Tối ưu hóa nguyên-hỗn hợp là nền tảng tính toán của các kỹ thuật được giới thiệu trong đồ án này.
Nhìn chung, một \textbf{Bài toán Quy hoạch Nguyên-Hỗn hợp (Mixed-Integer Program~-- MIP)} có thể rất khó
giải và thời gian chạy của bộ giải nguyên-hỗn hợp có thể tăng rất nhanh (theo cấp số nhân) cùng kích thước
bài toán.
Tuy nhiên, đây là trường hợp xấu nhất; trên thực tế, có thể xây dựng các MIP hiệu quả cao
được giải nhanh trong hầu hết các trường hợp mà chúng ta gặp phải.
Mục tiêu của phần này là hai mặt: thứ nhất, giới thiệu các kiến thức nền tảng về tối ưu hóa
nguyên-hỗn hợp (MIP là gì, cách phân loại MIP, và các thuật toán nào có thể dùng để giải chúng);
thứ hai, giới thiệu các ký hiệu và công cụ cần thiết để phân biệt giữa các MIP hiệu quả và không hiệu quả.

\subsection{Bài toán Quy hoạch Nguyên-Hỗn hợp}

Một MIP là bài toán tối ưu hóa với các biến liên tục $x \in \mathbb{R}^n$ và các biến rời rạc
(nguyên) $y \in \mathbb{Z}^m$.
Cho trước một hàm mục tiêu $f : \mathbb{R}^{n+m} \to \mathbb{R}$ và một tập ràng buộc
$\mathcal{S} \subseteq \mathbb{R}^{n+m}$, một MIP tổng quát có thể được phát biểu như sau:
\begin{subequations}
\label{eq:mip_full}
\begin{align}
\text{minimize} \quad & f(x, y) \label{eq:mip_full_obj}\\
\text{subject to} \quad & (x, y) \in \mathcal{S}, \label{eq:mip_full_con}\\
& y \in \mathbb{Z}^m. \label{eq:mip_full_int}
\end{align}
\end{subequations}
(Ràng buộc $y \in \mathbb{Z}^m$ được phát biểu tường minh vì, nếu không có quy định khác, các biến của
bài toán tối ưu luôn được giả định là có giá trị thực.)
Chúng ta gọi tập hợp
\[
\mathcal{T} := \mathcal{S} \cap \mathbb{R}^n \times \mathbb{Z}^m
\]
là \textit{tập chấp nhận được} (feasible set) của MIP~\eqref{eq:mip_full} và các phần tử
$(x, y) \in \mathcal{T}$ là các \textit{nghiệm chấp nhận được}.
Bài toán MIP là \textit{khả thi} nếu tồn tại nghiệm chấp nhận được và \textit{bất khả thi} trong
trường hợp ngược lại.
Một nghiệm chấp nhận được $(x^{\text{opt}}, y^{\text{opt}})$ là \textit{tối ưu} nếu
$f(x^{\text{opt}}, y^{\text{opt}}) \leq f(x, y)$ với mọi $(x, y) \in \mathcal{T}$.
Nếu nghiệm tối ưu tồn tại, ta ký hiệu $f^{\text{opt}} := f(x^{\text{opt}}, y^{\text{opt}})$
là \textit{giá trị tối ưu} của MIP.

Nếu bỏ ràng buộc~\eqref{eq:mip_full_int} khỏi MIP, ta thu được bài toán tối ưu sau:
\begin{subequations}
\label{eq:mip_relax}
\begin{align}
\text{minimize} \quad & f(x, y) \\
\text{subject to} \quad & (x, y) \in \mathcal{S}.
\end{align}
\end{subequations}
Bài toán này được gọi là \textbf{bài toán nới lỏng} (relaxation) của MIP.
Ta ký hiệu nghiệm tối ưu của bài toán nới lỏng là $(x^{\text{relax}}, y^{\text{relax}})$
và $f^{\text{relax}} := f(x^{\text{relax}}, y^{\text{relax}})$ là giá trị tối ưu của nó.
Lưu ý rằng
\[
f^{\text{relax}} \leq f^{\text{opt}},
\]
do việc bỏ một ràng buộc khỏi bài toán tối ưu (tối thiểu hóa) chỉ có thể làm giảm giá trị tối ưu.
Như được thảo luận ở phần dưới, độ chặt của bất đẳng thức này đóng vai trò trung tâm trong tính hiệu
quả của một MIP.
Nói một cách ngắn gọn, một MIP có thể được giải hiệu quả nếu bài toán nới lỏng của nó có thể được giải
nhanh và khoảng cách $f^{\text{opt}} - f^{\text{relax}} \geq 0$ là nhỏ.

Các MIP được phân loại theo các tính chất của hàm $f$ và tập $\mathcal{S}$,
những yếu tố này có thể ảnh hưởng đáng kể đến khả năng giải MIP một cách hiệu quả.
Dưới đây chúng ta định nghĩa các lớp MIP quan trọng nhất trong phạm vi đồ án này.

\subsection{Bài toán Quy hoạch Nhị phân-Hỗn hợp}

Một \textbf{Bài toán Quy hoạch Nhị phân-Hỗn hợp (Mixed-Boolean Program~-- MBP)} là một MIP
trong đó các biến rời rạc chỉ có thể nhận giá trị nhị phân: $y \in \{0, 1\}^m$.
Tương đương với đó, một MBP là bài toán có dạng~\eqref{eq:mip_full} với
\[
\mathcal{S} \subset \mathbb{R}^n \times [0, 1]^m.
\]
Tất cả các MIP mà chúng ta sẽ gặp trong đồ án này thực ra đều là các MBP.
Tuy nhiên, thuật ngữ MBP không phổ biến và, mặc dù về mặt kỹ thuật các bài toán của chúng ta là MBP,
chúng ta vẫn gọi chúng là MIP cho nhất quán với quy ước của lĩnh vực.

\subsection{Bài toán Quy hoạch Lồi Nguyên-Hỗn hợp}

Nếu hàm mục tiêu $f$ và tập ràng buộc $\mathcal{S}$ đều là lồi, ta gọi bài toán~\eqref{eq:mip_full}
là \textbf{Bài toán Quy hoạch Lồi Nguyên-Hỗn hợp (Mixed-Integer Convex Program~-- MICP)}.
MICP là một lớp cơ bản của MIP vì các bài toán nới lỏng của chúng là các bài toán tối ưu hóa lồi,
vốn trong hầu hết các trường hợp có thể được giải rất nhanh.
Để nhấn mạnh giả thiết tính lồi, ta gọi bài toán nới lỏng của một MICP là \textit{nới lỏng lồi}
(convex relaxation).
Sử dụng thuật toán Nhánh-và-Cận (Branch-and-Bound~-- BB), hầu hết các MICP có thể được giải một cách
đáng tin cậy đến mức tối ưu toàn cục, mặc dù thuật toán có thể mất thời gian đáng kể.

Trái ngược với lớp MICP, chúng ta gọi \textbf{Bài toán Quy hoạch Phi-lồi Nguyên-Hỗn hợp
(Mixed-Integer Non-Convex Program~-- MINCP)} là một MIP trong đó hàm mục tiêu $f$ và/hoặc tập ràng buộc
$\mathcal{S}$ là không lồi.
Hầu hết các MINCP đều không thể giải được trong thời gian thực tế; ngoại lệ chính là các MINCP
rất nhỏ trong đó tập chấp nhận được $\mathcal{T}$ có thể được rời rạc hóa mịn và tìm kiếm vét cạn.

\subsection{Bài toán Quy hoạch Nón Nguyên-Hỗn hợp}

Nếu hàm mục tiêu $f$ là tuyến tính và tập ràng buộc $\mathcal{S}$ là một tập lồi đóng dưới dạng nón,
bài toán~\eqref{eq:mip_full} được gọi là \textbf{Bài toán Quy hoạch Nón Nguyên-Hỗn hợp
(Mixed-Integer Conic Program~-- MIKP)}.
Các lớp con của nhóm bài toán này được phân loại như sau:

\begin{itemize}
    \item \textbf{Quy hoạch Tuyến tính Nguyên-Hỗn hợp (Mixed-Integer Linear Program~-- MILP)}
    khi $\mathcal{S}$ là một đa diện (polyhedron).
    \item \textbf{Quy hoạch Nón Bậc hai Nguyên-Hỗn hợp (Mixed-Integer Second-Order Cone Program~-- MISOCP)}
    khi $\mathcal{S}$ là tập lồi bậc hai.
    \item \textbf{Quy hoạch Bán xác định Nguyên-Hỗn hợp (Mixed-Integer Semidefinite Program~-- MISDP)}
    khi $\mathcal{S}$ là một spectrahedron.
\end{itemize}

Các bài toán nới lỏng của chúng được đặt tên theo cách tương tự; ví dụ, bài toán nới lỏng của một MILP
được gọi là \textit{nới lỏng tuyến tính} (linear relaxation).

MILP đóng vai trò đặc biệt quan trọng trong tối ưu hóa nguyên-hỗn hợp, quan trọng hơn cả vai trò
của LP trong tối ưu hóa lồi, vì hai lý do.
Lý do thứ nhất mang tính hình học: các đa diện là hình dạng đơn giản nhất có thể bao quanh một tập
hữu hạn các điểm, và do đó là ứng viên tự nhiên để mô hình hóa phần rời rạc của một MIP.
Lý do thứ hai mang tính thuật toán: thuật toán đơn hình (simplex) cho tối ưu tuyến tính có thể
được tích hợp vào thủ tục BB hiệu quả hơn so với thuật toán điểm trong được dùng cho các chương trình
nón tổng quát hơn.

Ngoài ra, một \textbf{Bài toán Quy hoạch Toàn phương Nguyên-Hỗn hợp (Mixed-Integer Quadratic
Program~-- MIQP)} là một MIP có dạng~\eqref{eq:mip_full} với $f$ là hàm bậc hai và $\mathcal{S}$ là
đa diện.
Các bài toán này có thể được biểu diễn dưới dạng MISOCP, nhưng tương tự như MILP, chúng có thể được
giải hiệu quả hơn bằng các phương pháp BB chuyên biệt khai thác tính đa diện của tập ràng buộc.
Ngày nay, các MISOCP và MISDP cũng có thể được giải khá hiệu quả ngay cả với các bộ giải mã nguồn mở
(xem, ví dụ, Pajarito~\cite{lubin2018}).

